
\chapter{Lean Source Code}
\section{Auxiliary Lean Proofs}
This section contains all proofs needed to mechanize the definitions and their surrounding semantic functions.


\subsection*{State}
\label{app:lean:state}
State is loosely defined in the textbook as one of multiple different choices,
In this thesis we opted to use the Functional approach as Lean is very good at handling it.
To do this we need to introduce some structure, first of which is the definition of what \emph{State} is.

\begin{Lean_Code}
def State := Var @\to@ @ \(\mathbb{Z}\)@
\end{Lean_Code}
Now that we have the functional definition, it can be helpful to create a "standard implementation", in this case a state which maps all values to \(0\):

\begin{Lean_Code}
def default_state : State :=
  fun _ => 0
\end{Lean_Code}
By adding a default state that satisfies the total function property we can then use the assignment operator to create the states used in specific examples.

In this case the assignment was defined as follows:

\begin{Lean_Code}
def assign (s : State) (x : Var) (z : @\(\mathbb{Z}\)@) : State :=
  fun (v: Var) => if v = x then z else s v
\end{Lean_Code}
with the following notation:

\begin{Lean_Code}
notation:max s "[" x "@\mapsto@" z "]" => assign s x z
\end{Lean_Code}
Next we need some scaffolding proofs to be able to reason if 2 states are equivalent or not.
This is provided by the two theorems:
\begin{Lean_Code}
theorem assign_shadow (s : State) (x : Var) (z1 z2 : @\(\mathbb{Z}\)@) :
  s[x @\mapsto@ z1][x @\mapsto@ z2] = s[x @\mapsto@ z2] := by @\dots@
  apply funext
  intro v
  repeat rw [assign]
  split_ifs
  · rfl
  · rfl
\end{Lean_Code}
and
\begin{Lean_Code}
theorem assign_comm (s : State) (x y : Var) @\leanbreak@ (z1 z2 : @\(\mathbb{Z}\)@) (h : x @\neq@ y) :
  s[x @\mapsto@ z1][y @\mapsto@ z2] = s[y @\mapsto@ z2][x @\mapsto@ z1] := by @\dots@
  apply funext
  intro v
  repeat rw [assign]
  split_ifs with h1 h2
  -- h1 : v = y
  -- h2 : v = x
  · show z2 = z1
    subst h1 h2
    contradiction
  · show z2 = z2
    rfl
  · show z1 = z1
    rfl
  · show s v = s v
    rfl
\end{Lean_Code}
which together allows us to know that the following is true:
\begin{align*}
  s[x \mapsto 0][y \mapsto 6] &= s[x \mapsto 3][y \mapsto 1][x \mapsto 2][y \mapsto 3][x \mapsto 1][y \mapsto 6][x \mapsto 0][y \mapsto 6] \\
  \text{assign\_comm}   &= s[x \mapsto 3][x \mapsto 2][y \mapsto 1][y \mapsto 3][x \mapsto 1][y \mapsto 6][x \mapsto 0][y \mapsto 6] \\
  \text{assign\_shadow} &= s[x \mapsto 2][y \mapsto 1][y \mapsto 3][x \mapsto 1][y \mapsto 6][x \mapsto 0][y \mapsto 6] \\
  \text{assign\_shadow} &= s[x \mapsto 2][y \mapsto 3][x \mapsto 1][y \mapsto 6][x \mapsto 0][y \mapsto 6] \\
  \text{assign\_comm}   &= s[x \mapsto 2][x \mapsto 1][y \mapsto 3][y \mapsto 6][x \mapsto 0][y \mapsto 6] \\
  \text{assign\_shadow} &= s[x \mapsto 1][y \mapsto 3][y \mapsto 6][x \mapsto 0][y \mapsto 6] \\
  \text{assign\_shadow} &= s[x \mapsto 1][y \mapsto 6][x \mapsto 0][y \mapsto 6] \\
  \text{assign\_comm}   &= s[x \mapsto 1][x \mapsto 0][y \mapsto 6][y \mapsto 6] \\
  \text{assign\_shadow} &= s[x \mapsto 0][y \mapsto 6][y \mapsto 6] \\
  \text{assign\_shadow} &= s[x \mapsto 0][y \mapsto 6]
\end{align*}

\section{Exercise 2.21}
\begin{Lean_Code}
/-!
Exercise 2.21 (Essential)
Prove that if `@\(\langle\)@S@\(_1\)@, s@\(\rangle\)@ @\(\to_{sos}\)@[k] s'` then `@\(\langle\)@S@\(_1\)@; S@\(_2\)@, s@\(\rangle\)@ @\(\to_{sos}\)@[k] @\(\langle\)@S@\(_2\)@, s'@\(\rangle\)@`.
The execution of `S@\(_1\)@` is not influenced by the statement following it.
-/
lemma seq_exec_preserve_right {S@\(_1\)@ S@\(_2\)@ : Stmt} {s s' : State} {k : Nat}
  (h : @\(\langle\)@S@\(_1\)@, s@\(\rangle\)@ @\(\to_{sos}\)@[k] s') :
  @\(\langle\)@S@\(_1\)@ ; S@\(_2\)@, s@\(\rangle\)@ @\(\to_{sos}\)@[k] @\(\langle\)@S@\(_2\)@, s'@\(\rangle\)@ := by
  induction k generalizing s S@\(_1\)@ with
  | zero =>
    -- @\(\langle\)@S@\(_1\)@, s@\(\rangle\)@ @\(\to_{sos}\)@[0] s' is impossible because Step ≠ Final
    cases h
  | succ k ih =>
    -- We have @\(\langle\)@S@\(_1\)@, s@\(\rangle\)@ @\(\to_{sos}\)@ @\gamma@' @\(\to_{sos}\)@[k] s'
    cases h with
    | step h_step h_rest =>
      rename_i @\gamma@'
      -- Now we branch on whether @\gamma@' is a final state or a continuing step
      cases @\gamma@' with
      | final s_next =>
        -- Case 1: S@\(_1\)@ finishes in one step (@\gamma@' = s_next)
        -- This matches the [comp2] rule: @\(\langle\)@S@\(_1\)@; S@\(_2\)@, s@\(\rangle\)@ @\(\to_{sos}\)@ @\(\langle\)@S@\(_2\)@, s_next@\(\rangle\)@
        -- Since h_rest is @\(\langle\)@s_next, ...@\(\rangle\)@ @\(\to_{sos}\)@[k] s', and s_next is final,
        -- k must be 0 and s_next = s'.
        cases h_rest
        exact small_step_k.step (small_step.comp2 h_step) small_step_k.refl
      | step S_next s_next =>
        -- Case 2: S@\(_1\)@ takes one step to a new configuration @\(\langle\)@S_next, s_next@\(\rangle\)@
        -- This matches the [comp1] rule: @\(\langle\)@S@\(_1\)@; S@\(_2\)@, s@\(\rangle\)@ @\(\to_{sos}\)@ @\(\langle\)@S_next; S@\(_2\)@, s_next@\(\rangle\)@
        apply small_step_k.step
        · exact small_step.comp1 h_step
        · -- Now we need: @\(\langle\)@S_next; S@\(_2\)@, s_next@\(\rangle\)@ @\(\to_{sos}\)@[k] @\(\langle\)@S@\(_2\)@, s'@\(\rangle\)@
          -- This is exactly our IH applied to the remaining steps h_rest
          exact ih h_rest

end exercise_2_21

lemma small_step_k_to_star {@\gamma@@\(_1\)@ @\gamma@@\(_2\)@ : Config} {k : Nat}
  (h : small_step_k @\gamma@@\(_1\)@ k @\gamma@@\(_2\)@) : small_step_star @\gamma@@\(_1\)@ @\gamma@@\(_2\)@ := by
  induction h with
  | refl => exact small_step_star.refl
  | step s h_rest ih => exact small_step_star.step s ih


lemma star_to_small_step_k {@\gamma@@\(_1\)@ @\gamma@@\(_2\)@ : Config}
  (h : small_step_star @\gamma@@\(_1\)@ @\gamma@@\(_2\)@) : @\exists@ k, small_step_k @\gamma@@\(_1\)@ k @\gamma@@\(_2\)@ := by
  induction h with
  | refl =>
    exists 0
    exact small_step_k.refl
  | step s h_rest ih =>
    let @\(\langle\)@k, hk@\(\rangle\)@ := ih
    exists k + 1
    exact small_step_k.step s hk

lemma small_step_star_trans {@\gamma@@\(_1\)@ @\gamma@@\(_2\)@ @\gamma@@\(_3\)@ : Config}
  (h_left : small_step_star @\gamma@@\(_1\)@ @\gamma@@\(_2\)@) (h_right : small_step_star @\gamma@@\(_2\)@ @\gamma@@\(_3\)@) :
  small_step_star @\gamma@@\(_1\)@ @\gamma@@\(_3\)@ := by
  induction h_left with
  | refl => exact h_right
  | step s h_rest ih => exact small_step_star.step s (ih h_right)

lemma exercise_2_21 {S@\(_1\)@ S@\(_2\)@ : Stmt} {s s' : State}
  (h : @\(\langle\)@S@\(_1\)@, s@\(\rangle\)@ @\(\to_{sos}\)@* s') : @\(\langle\)@S@\(_1\)@ ; S@\(_2\)@, s@\(\rangle\)@ @\(\to_{sos}\)@* @\(\langle\)@S@\(_2\)@, s'@\(\rangle\)@ := by
  -- 1. Convert star to k-step
  let @\(\langle\)@k, hk@\(\rangle\)@ := star_to_small_step_k h
  -- 2. Explicitly tell Lean to use S@\(_2\)@ from the goal
  -- We use (S@\(_2\)@ := S@\(_2\)@) to map the lemma's parameter to our local variable
  let h_k := seq_exec_preserve_right (S@\(_2\)@ := S@\(_2\)@) hk
  -- 3. Convert back to star
  exact small_step_k_to_star h_k
\end{Lean_Code}